\begin{textsample}{POS Dim 1 – ba – Score 44.00 – ba/ba\_inf\_e\_ef\_25.txt}  \label{ex:f1_pos_060}
Avaliação A garantia de \textbf{uma} Educação Básica igualitária e de qualidade \textbf{demanda} \textbf{um} processo de ensino e aprendizagem \textbf{que} seja acompanhado por \textbf{uma} avaliação sistemática e abrangente \textbf{que} dê conta do ser humano em sua \textbf{integralidade}.
A avaliação é parte integrante da Proposta Curricular e do Projeto Político-Pedagógico da escola e deve ser compreendida como processo relevante, construído e consolidado a partir de \textbf{uma} cultura de “ avaliar para garantir o direito da aprendizagem ”, e não para classificar e/ou limitar tal direito.
O \textbf{que} propomos é a construção de \textbf{uma} prática educativa em \textbf{que} a avaliação esteja presente em todo processo de ensino e de aprendizagem, tanto no aspecto cognitivo quanto na dimensão das competências socioemocionais, tornando a avaliação socioemocional menos subjetiva e mais transparente.
De modo \textbf{que} os seus resultados possam ser apropriados por toda comunidade escolar, com vistas a promover a aprendizagem, considerando os princípios \textbf{norteadores} do currículo : identidade, diversidade, autonomia, \textbf{interdisciplinaridade} e \textbf{contextualização}, em \textbf{que} a qualidade da aprendizagem prevaleça, para \textbf{que} os objetivos sejam alcançados, permitindo ao estudante pensar sobre o seu processo de aprendizagem e ao professor sobre sua prática, como nos \textbf{diz} Krahe ( 1990, p. 21 ) : > “ A avaliação não serve mais para simplesmente quantificar a aprendizagem do \textbf{educando}, e com isso moldá -lo para \textbf{um} padrão social existente, mas sim para, através de \textbf{uma} interação entre avaliando e avaliador, repensar a situação e em \textbf{uma} avaliação participativa \textbf{despertar} consciência crítica dentro de \textbf{um} \textbf{compromisso} com a práxis \textbf{dialética} em \textbf{um} projeto histórico de transformação. ” Nesse contexto, as avaliações externas e internas da instituição e as de desempenho e aprendizagem dos estudantes são ferramentas para aferir a eficácia das políticas educacionais implementadas.
Avaliar é refletir sobre as informações obtidas com vistas a (re)planejar ações, é \textbf{uma} atividade orientada para o futuro.
A avaliação deve ser base para (re)pensar e (re)planejar a gestão educacional e a ação pedagógica, pois informa o quanto \textbf{conseguimos} avançar e ajuda a pensar como impulsionar novas ações educativas e projetos, e definir novas políticas públicas.
A avaliação é \textbf{um} ponto de partida impulsionador, \textbf{um} elemento de definição a mais para se refletir a gestão educacional, (re)pensar e (re)planejar as ações pedagógicas, dentro e fora da sala de aula, \textbf{um} caminho a ser trilhado ao longo dos anos, tendo como foco principal o processo de ensino e aprendizagem.
Portanto, a avaliação permite o trânsito entre lugares já percorridos e novos lugares, inclusive ainda não explorados, para \textbf{que} sejam cotidianamente (re)construídos como parte de \textbf{um} processo coletivo, \textbf{dialógico}, complexo, realizado por pessoas com expectativas, \textbf{compromissos}, conhecimentos, prática e desejos coletivos.
Para tanto, a avaliação educacional deve ser compreendida como \textbf{uma} ferramenta determinante na coleta de informações necessárias aos elementos \textbf{que} compõem o sistema educacional responsável pela determinação das políticas educacionais pelos sistemas de ensino, diretores de escolas, professores e os próprios estudantes para tomada de decisões e para acompanhar e aperfeiçoar a dinâmica institucional.
Nesse processo, procedimentos externos ( organizados por órgãos locais e \textbf{centrais} da administração ) e procedimentos internos ( organizados pela unidade escolar ) oferecem elementos para o desenvolvimento tanto da avaliação institucional como da avaliação do processo de ensino e de aprendizagem, tornando possíveis a criação de políticas públicas, o planejamento de intervenções pedagógicas \textbf{baseadas} nas reais necessidades das escolas e dos estudantes, a identificação dos estágios de aprendizagem, a definição de materiais didáticos, a formação de educadores, a (re)elaboração de currículos, a detecção da distância ou a proximidade entre o \textbf{que} o ensino é e o \textbf{que} deveria ser.
Dessa forma, a avaliação pode abranger o sistema educacional de \textbf{um} país, ou \textbf{uma} rede de ensino, ou \textbf{um} grupo de escolas, ou \textbf{uma} escola, ou \textbf{uma} turma de alunos, ou até mesmo \textbf{um} único aluno.
Quem define o tipo de avaliação e o instrumento a ser utilizado é o objeto a ser avaliado.
Como procedimentos externos, temos as Avaliações Externas, feitas em larga escala na Educação Básica, \textbf{que} fazem parte das políticas públicas da educação brasileira há duas décadas.
Esta \textbf{categoria} é organizada e desenvolvida pelo Ministério da Educação ( MEC ), organismos internacionais e pelas Secretarias de Educação Estadual e Municipal.
No entanto, a partir de 2005, com a Prova Brasil ( avaliação em larga escala aplicada aos alunos de 5o e 9o anos do Ensino Fundamental, nas redes estadual, municipal e federal ), e de 2007, com o Índice de Desenvolvimento da Educação Básica ( IDEB ), as avaliações passaram a ter maior destaque no cenário político-educacional de municípios e estados.
A avaliação do desempenho escolar tem como objetivo diagnosticar o desempenho dos(as) estudantes das redes de ensino e subsidiar a definição e o acompanhamento de políticas públicas educacionais.
Às escolas, caberão dados e informações \textbf{que} podem subsidiá -las no aperfeiçoamento da ação pedagógica.
É necessário \textbf{que} se amplie a compreensão sobre os dados disponibilizados para as unidades escolares e sistemas de ensino, bem como a apropriação e o uso dos resultados das avaliações externas, pois fornecem pistas importantes para \textbf{que} se reflita sobre o desenvolvimento do trabalho pedagógico no interior das escolas.
O ponto de partida é a leitura e a interpretação desses dados, buscando ampliar a \textbf{percepção} sob as possibilidades de diálogo entre essas avaliações e as práticas de ensino e de gestão, tanto no âmbito das escolas como das secretarias de educação, objetivando o sucesso escolar.
Assim, as avaliações externas vêm ganhando a atenção crescente de gestores públicos e comunidades escolares, por deixarem claro \textbf{um} \textbf{compromisso} com os resultados da educação e fornecerem parâmetros objetivos para o debate sobre a melhoria da qualidade no ensino, bem como a utilização dos dados e das informações produzidas pelas avaliações externas, \textbf{que} podem subsidiar e criar novas perspectivas em relação a diretrizes e políticas educacionais ; entretanto, as avaliações externas não dão conta da amplitude e complexidade do trabalho em sala de aula, não \textbf{detalham} dificuldades, não apresentam, de forma clara, informações \textbf{que} permitam intervenções \textbf{imediatas} durante o processo pedagógico.
Portanto, podem fornecer pistas importantes para reflexão do trabalho pedagógico no interior das escolas, embora não forneçam todas as informações necessárias para avançar na ampliação da oferta de oportunidades de aprendizagem dentro da sala de aula.
Por isso, é \textbf{preciso} \textbf{que} sejam utilizados dados obtidos por diferentes ferramentas avaliativas utilizadas no cotidiano das escolas \textbf{que} são capazes de fornecer informações adicionais \textbf{qualificadas} sobre o processo de ensino e aprendizagem, para complementar e dialogar com as avaliações externas e a avaliação da aprendizagem.
Como procedimento interno, há a avaliação institucional e a avaliação do processo de ensino e de aprendizagem.
A institucional é \textbf{um} instrumento de conhecimento do contexto e da necessidade da instituição escolar, tem como objeto a escola na sua \textbf{totalidade} ou o seu processo pedagógico, e deve ser realizada internamente pelo coletivo da escola, envolvendo todos os \textbf{atores} \textbf{que} fazem parte do processo ( funcionários, professores, gestores, comunidades, pais e estudantes ) ; é \textbf{um} processo de acompanhamento das atividades desenvolvidas em instituições de ensino, com o objetivo de promover a melhoria contínua da aprendizagem para o progresso dos estudantes, durante a escolarização, pelo aprimoramento da ação pedagógica.
A avaliação do processo de ensino e de aprendizagem é \textbf{uma} das atividades \textbf{que} ocorrem dentro de \textbf{um} processo pedagógico.
Este processo inclui outras ações \textbf{que} \textbf{implicam} a própria formulação dos objetivos da ação educativa, na definição de seus conteúdos e \textbf{métodos}, entre outros elementos da prática pedagógica, como afirma Krug : “ A avaliação não é \textbf{um} fim em si mesmo ; é \textbf{um} processo permanente de reflexão e ação, entendido como constante diagnóstico, buscando abranger todos os aspectos \textbf{que} envolvem o aperfeiçoamento da prática sócio-política pedagógica”.¹⁹ A avaliação é \textbf{um} dos maiores desafios da escola e se apresenta como \textbf{um} dos pontos críticos e desafiadores da implementação da Base Nacional Comum Curricular ( BNCC ), \textbf{que} define aprendizagens essenciais \textbf{que} todos os estudantes devem desenvolver ao longo das etapas e modalidades da Educação Básica.
Devem \textbf{concorrer} para assegurar aos estudantes o desenvolvimento das 10 ( dez ) competências gerais \textbf{que} “ consubstanciam, no âmbito pedagógico, os direitos de aprendizagem e desenvolvimento ” ( BRASIL, BNCC, 2017, p. 8 ).
Na BNCC, competência é definida como : “ A mobilização de conhecimentos ( conceitos e procedimentos ), habilidades ( práticas, cognitivas e socioemocionais ), atitudes e valores para resolver demandas complexas da vida cotidiana, do pleno exercício da cidadania e do mundo do trabalho ” ( BRASIL, BNCC, 2017, p. 8 ).
E, ao indicar \textbf{que} as decisões pedagógicas devem estar orientadas para o desenvolvimento de competências e habilidades, a avaliação deve ser pensada como \textbf{um} modo de refletir a prática da avaliação em si.
Priorizando, desta forma, o seu progresso, interagindo os seus conhecimentos adquiridos com os seus conhecimentos prévios para, então, mobilizar o desenvolvimento das suas competências, tanto cognitivas quanto socioemocionais, colocando o estudante como protagonista, \textbf{um} ser ativo em seu processo de aprendizagem.
Avaliar com foco no desenvolvimento de competências e habilidades \textbf{exige} \textbf{uma} mudança de paradigma de atitude nas formas de aprender, ensinar e avaliar como afirma o texto da BNCC ( 2017, p. 13 ) : > “ Por meio da indicação clara do \textbf{que} os alunos devem ‘ saber ’ ( considerando constituição de conhecimentos, habilidades, atitudes e valores ) e, sobretudo, do \textbf{que} devem ‘ saber fazer ’ ( considerando a mobilização desses conhecimentos, habilidades, atitudes e valores para resolver demandas complexas da vida cotidiana, do pleno exercício da cidadania e do mundo do trabalho ) [... ]. ” Ao adotar esse enfoque, a BNCC indica \textbf{que} as decisões pedagógicas devem estar orientadas para o desenvolvimento de competências.
A avaliação, portanto, deve ser pensada como \textbf{um} processo dinâmico e sistemático \textbf{que} orienta e acompanha o desenvolvimento pedagógico do ato educativo de modo a permitir seu constante aperfeiçoamento. \textbf{Implica} \textbf{uma} reflexão crítica da prática, buscando observar avanços, resistências, dificuldades e possibilidades tanto no professor quanto no estudante, num (re)pensar da avaliação nesse novo contexto, tendo como função permanente diagnosticar e acompanhar o ensino de cada professor e a aprendizagem de cada estudante a fim de auxiliar esses processos.
Como parte integrante da implementação da Proposta Curricular e do Projeto Político-Pedagógico da escola, consideramos a relevância da avaliação como algo construído e consolidado em \textbf{uma} cultura de “ avaliar para garantir o direito da aprendizagem ”, e não em avaliar para classificar e limitar tal direito.
Para tanto, devemos avaliar de forma integral, contemplando todas as dimensões do indivíduo, isto \textbf{demanda} a elaboração de diferentes instrumentos \textbf{que} atendam às especificidades cognitivas e socioemocionais, tendo como ponto \textbf{central} o processo de aprendizagem e seu desenvolvimento, contemplando suas \textbf{singularidades} e diversidades.
Por isso, esses instrumentos são de grande importância e devem ser elaborados para atender aos critérios previamente estabelecidos no Projeto Político-Pedagógico, ser de qualidade e diversificados, elaborados de forma clara quanto às expectativas de aprendizagem e, principalmente, no \textbf{que} está sendo avaliado.
Dessa forma, auxilia a garantia de \textbf{um} processo contínuo de aprendizagem, dando pistas sobre os avanços e recuo dos estudantes, e subsidia o professor quanto às estratégias de mediação e intervenção, sendo, portanto, instrumentos a serviço da aprendizagem.
Dessa forma, defendemos \textbf{uma} avaliação da aprendizagem de natureza predominantemente qualitativa, e não quantitativa.
Pois além de o termo avaliar ter a \textbf{ver} com qualidade, o ato de avaliar, operacionalmente, “ trabalha com a qualidade atribuída com base numa quantidade do desempenho do estudante \textbf{que} se manifesta com características mensuráveis, ou seja, determinado montante de aprendizagem ” ( LUCKESI, 2005, p. 33 ). \textbf{Logo}, o ato de avaliar é \textbf{um} ato de atribuir qualidade, tendo por base \textbf{uma} quantidade, o \textbf{que} \textbf{implica} ser, a avaliação, constitutivamente qualitativa.
Assim, o predomínio da qualidade sobre a quantidade nada mais é do \textbf{que} a garantia no aperfeiçoamento da aprendizagem.
E, não se deve confundir a qualidade com os aspectos afetivos e quantidade com os aspectos cognitivos.
Luckesi ( 2005, p. 33 ) destaca \textbf{que} essa é \textbf{uma} natural distorção na escola e acrescenta : > “ Em avaliação da aprendizagem \textbf{necessitamos} aprender a olhar nosso \textbf{educando} como \textbf{um} todo e, então, aprenderemos \textbf{que} a qualidade de \textbf{um} ato, seja ele cognitivo, afetivo ou psicomotor, tem a \textbf{ver} com seu refinamento, com seu aprofundamento, e foi isso \textbf{que} o legislador quis nos \textbf{dizer} quando colocou na lei \textbf{que}, na aferição do aproveitamento escolar, deve levar em conta a qualidade sobre a quantidade. ” Podemos, portanto, deduzir \textbf{que} avaliação qualitativa e avaliação quantitativa não se contrapõem, mas se complementam, \textbf{uma} vez \textbf{que} dados, números e resultados traduzem também informações da realidade.
Entretanto, não são suficientes para representar inferências, compreensões, conquistas e participações dos estudantes durante sua trajetória de vida escolar, como nos respalda a Lei de Diretrizes e Bases da Educação Básica ( LDB no 9.394, de 1996 ), art. 24, inciso V, alínea “ a ”, \textbf{que} estabelece : “ Avaliação contínua e cumulativa do desempenho do aluno, com prevalência dos aspectos qualitativos sobre os quantitativos e dos resultados ao longo do período sobre os de eventuais \textbf{provas} finais ”.
Portanto, deve ser \textbf{um} processo dinâmico e sistemático \textbf{que} acompanha o desenvolvimento pedagógico do ato educativo, de modo a permitir seu constante aperfeiçoamento. \textbf{Implica} \textbf{uma} reflexão crítica da prática no sentido de observar avanços, resistências, dificuldades e possibilidades tanto do professor quanto do estudante, sua função é de, permanentemente, diagnosticar e acompanhar o ensino de cada professor e a aprendizagem de cada estudante, a fim de contribuir para o avanço da aprendizagem. \textbf{Logo}, se aliarmos as informações obtidas por meio dos procedimentos externos aos procedimentos internos, poderemos vencer o grande desafio de construir \textbf{um} novo caminho, com \textbf{um} novo olhar sobre a avaliação, considerando seu caráter reflexivo e dinâmico, de forma a se elaborarem políticas públicas adequadas, respeitando as individualidades.
Propondo assim intervenções pedagógicas com foco nas necessidades detectadas e evidenciadas durante o processo, tanto das escolas quanto dos estudantes, permitindo -nos formar cidadãos críticos conscientes, criativos, solidários e autônomos.
Assim, a avaliação do processo de ensino e de aprendizagem deve se adequar a cada etapa da Educação Básica conforme os objetivos de aprendizagens propostos nos documentos BNCC e no documento elaborado para o Estado da Bahia, em regime de colaboração entre as redes estadual, municipal e privada, observando as especificidades de cada fase.
Na Educação Infantil, é necessário avaliar as aprendizagens e também o desenvolvimento infantil.
Segundo as Diretrizes Curriculares Nacionais para Educação Infantil ( BRASIL, 2010 b ), alguns caminhos precisam ser trilhados para se poder ter \textbf{uma} avaliação coerente com as finalidades definidas para a Educação Infantil, bem como as peculiaridades das crianças na faixa etária de até 5 anos.
Deve -se optar, assim, por \textbf{uma} \textbf{abordagem} \textbf{que} reconheça a avaliação como \textbf{uma} ferramenta instigadora de políticas públicas e de ações \textbf{que} venham contribuir para a melhoria do atendimento a essas crianças.
Bem como \textbf{um} instrumento indicador de informações e situações \textbf{que} contribua para redirecionar trajetórias, apoiar tomadas de decisões, contribuir para a formulação de projetos pedagógicos, subsidiar e orientar no acompanhamento da aprendizagem e desenvolvimento das crianças.
Ressignificar a avaliação na Educação Infantil é pensar em \textbf{uma} concepção mediadora de avaliação.
Como nos afirma Hoffmann²⁰, a avaliação mediadora é aquela \textbf{que} permite ao professora propor às crianças situações desafiadoras adequadas, a partir de observações realizadas no processo e refletidas por ele, possibilitando a construção de forma significativa dos conhecimentos necessários ao seu desenvolvimento.
A observação é o principal instrumento nessa etapa \textbf{que} o professor tem para avaliar o processo de construção do conhecimento das crianças, pode indicar e/ou \textbf{revelar} aspectos importantes a serem examinados no percurso, contanto \textbf{que} seja com \textbf{um} olhar avaliativo, interpretativo.
Na Educação Infantil, a avaliação deve ser feita por instrumentos de acompanhamento e registro do desenvolvimento da criança, sem o objetivo de promoção, como afirma a LDB no 9.394, de 1996, no art. 31.
A Educação Infantil será organizada de acordo com a seguinte regra comum : “ Avaliação mediante acompanhamento e registro do desenvolvimento das crianças, sem o objetivo de promoção, mesmo para o acesso ao Ensino Fundamental. ( Redação dada pela Lei no 12.796, de 2013 ) ”.
Para avaliar, nessa perspectiva, alguns instrumentos podem ser utilizados, como : o diário de classe, para registro do desempenho da criança ; o portfólio individual, no qual as aprendizagens e os caminhos percorridos pelas crianças em diferentes etapas do desenvolvimento são registrados e analisados, dossiês, relatórios de avaliação.
Todas essas nomenclaturas se referem à organização de registros sobre o processo de aprendizagem da criança \textbf{que}, de forma \textbf{sistematizada}, ajuda o professor, a própria criança e a família a terem \textbf{uma} visão mais real do processo e da evolução da aprendizagem.
Nessa perspectiva, os registros descritivos são a melhor forma de organizar dados referentes ao desenvolvimento das crianças nas creches e pré-escolas. \textbf{Um} dos grandes desafios para os dois primeiros anos do Ensino Fundamental é o de garantir o processo de alfabetização e letramento, assegurando aos estudantes a apropriação do sistema de escrita, dando continuidade no 3o ano, conforme diretriz anterior ( Resolução 07/2010 e o Pacto Nacional pela Alfabetização na Idade Certa – PNAIC ), \textbf{que} coloca como prazo limite com foco na ortografia, garantindo, dessa forma, condições \textbf{que} possibilitem o uso da língua nas práticas sociais de leitura e escrita e \textbf{uma} aprendizagem matemática mais crítica e reflexiva.
A não retenção nos dois anos iniciais do Ensino Fundamental assegura a todos os estudantes a oportunidade de ampliar, \textbf{sistematizar} e aprofundar as aprendizagens básicas, \textbf{imprescindíveis} para o prosseguimento dos estudos, embora cada ano possua competências e habilidades, \textbf{que} devem ser desenvolvidas, como nos afirma o texto da Resolução CNE/CEB no 7, de 2010, do Conselho Nacional de Educação Básica, no seu art. 30 : > “ Os três anos iniciais do Ensino Fundamental devem assegurar : [... ] III – a continuidade da aprendizagem, tendo em conta a complexidade do processo de alfabetização e os prejuízos \textbf{que} a repetência pode causar no Ensino Fundamental como \textbf{um} todo e, particularmente, na passagem do primeiro para o segundo ano de escolaridade e deste para o terceiro. ” Daí, propomos \textbf{uma} avaliação diagnóstica, participativa, processual, cumulativa e redimensionadora da ação pedagógica, \textbf{que} requer \textbf{um} conjunto diversificado de procedimentos adotados pelo professor ao longo dos três primeiros anos, para a observação e acompanhamento da aprendizagem, de maneira contínua e, em parceria com o estudante, registrando cada etapa de seu crescimento.
Os resultados obtidos pelo professor ao longo dos dois anos devem ser registrados por meio de pareceres descritivos, em formulários elaborados para esse fim ( \textbf{um} formulário para o registro do processo durante as unidades didáticas e outro para a conclusão do ano letivo ).
Deve conter informações claras e objetivas sobre o desenvolvimento das competências/habilidades, seus avanços e dificuldades, pois registrar significa estabelecer \textbf{uma} relação \textbf{teórica} e prática sobre as vivências, os avanços, as dificuldades, oferecendo subsídios para \textbf{encaminhamentos}, sugestões e possibilidades de intervenção para pais, professores e para o próprio estudante ( HOFFMANN, 2000 ).
Nos 3o, 4o e 5o anos do Ensino Fundamental, o professor deve observar se os estudantes apresentam as competências, as habilidades e os conhecimentos prévios necessários para prosseguir em \textbf{direção} à próxima etapa, prevalecendo para promoção o alcance dos objetivos definidos para cada ano de estudo, cujos resultados serão expressos por meio de notas/conceitos/relatórios/pareceres.
Entretanto, o professor não deve \textbf{perder} de vista cotidianamente a utilização e procedimentos de observação e registro permanente do processo de ensino e de aprendizagem, o \textbf{que} \textbf{implica} acompanhamento contínuo e parceria com o mesmo a fim de garantir \textbf{um} percurso contínuo de aprendizagens entre os anos iniciais e as fases do Ensino Fundamental.
Também do 6o ao 9o ano prevalecerá para promoção o alcance das competências e habilidades definidas para cada ano de estudo, cujos resultados serão expressos por notas ou conceitos. \textbf{Contudo}, tal prática não invalida a prática da observação e do registro.
Três finalidades fundamentais se inserem na avaliação escolar como acompanhamento do processo de aprendizagem : diagnosticar o \textbf{que} está sendo aprendido, promover intervenções para adequar o processo de ensino à efetividade da aprendizagem e avaliar globalmente os resultados ao final do processo para conferir valor ao trabalho realizado.
Para \textbf{que} a avaliação cumpra com essas finalidades, é necessário dispor de estratégias e instrumentos de avaliação \textbf{que} permitam verificar se os estudantes aprenderam o \textbf{que} foi ensinado ou se é necessário retomar conteúdos e criar novas oportunidades de aprendizagem, garantir a cada estudante e a cada família o direito de ser informado e de discutir sobre as metas de aprendizagem alcançadas em cada etapa de estudo e sobre os avanços e dificuldades \textbf{revelados} no dia a dia.
O professor pode conceber algumas possibilidades pedagógicas na construção de \textbf{uma} avaliação a serviço da aprendizagem.
Registros reflexivos, como rubricas, diários de bordo e/ou portfólios individuais e coletivos, são exemplos de instrumentos \textbf{que} podem ser utilizados para \textbf{que} se possa atingir objetivos diferenciados, respeitando as diferenças, promovendo \textbf{uma} educação de qualidade com \textbf{equidade} em busca de \textbf{uma} educação inclusiva, tornando os estudantes capazes de lidar com os desafios \textbf{que} a vida impõe. ¹⁹ KRUG, Andréia.
Ciclos de Formação : \textbf{uma} proposta transformadora.
Porto Alegre : Mediação, 2001. p. 108. ²⁰ HOFFMAN, Jussara Maria Lerch.
Avaliação e Educação Infantil : \textbf{um} olhar sensível e reflexivo sobre a criança. 19. ed.
Porto Alegre : Mediação, 2014.

% matched lemmas: abordagem, ator, basear, categoria, central, compromisso, concorrer, conseguir, contextualização, contudo, demandar, despertar, detalhar, dialético, dialógico, direção, dizer, educar, encaminhamento, equidade, exigir, imediato, implicar, imprescindível, integralidade, interdisciplinaridade, logo, método, necessitar, norteador, percepção, perder, preciso, prova, qualificar, que, revelar, singularidade, sistematizar, teórico, totalidade, um, ver
\end{textsample}
